% Converted from README.md
\subsection{C++ primer - kalkulator}\label{c-primer---kalkulator}

Naš kod za kalkulator je nepodesan za jedinično testiranje jer se mnogo posla zapravo obavlja u main funkciji. Da bismo sve delove testirali moramo kod izdeliti na funkcije koje obavljaju po jednu celinu. Već postoje funkcije za svaku pojedinačnu operaciju koje imaju dva argumenta, njih nećemo dirati.

U \texttt{main} funkciji test podaci se učitavaju sa standardnog ulaza i sve se ispisuje na standardni izlaz. Da bismo mogli u jediničnim testovima da zadajemo svoje ulaze, promenićemo da funkcije ne učitavaju sa \texttt{std::cin} već sa proizvoljnog toka tipa \texttt{std::istream}, i da ispisuju ne na \texttt{std::cout}, već na proizvoljni tok tipa \texttt{std::ostream}. Na taj način ćemo izbeći učitavanje sa standardnog ulaza, a moći ćemo da unapred zadamo ulaz na kom se testira.

Iz tog razloga već postojećoj funkciji \texttt{showChoices} treba dodati argument \texttt{std::ostream\ \&\ ostr} na koji će ispisivati ponudene opcije umesto na \texttt{std::cout}. Vidimo da se u \texttt{main} funkciji nakon prikaza opcije očekuje unos izbora operacije i proverava ispravosti unetog podatka. Taj deo objedinjen izdvajamo u posebnu funkciju \texttt{readChoice} koja će pozivati funkciju \texttt{showChoices} i učitavati opciju sa ulaznog toka sve dok se ne unese jedna od ispravnih cifara.

Nakon toga bi trebalo da se unesu dva operanda nad kojima će se primeniti izabrana operacija. To se može izdvojiti u posebnu funkciju \texttt{readOperands} koja će učitati dva realna broja sa ulaznog toka. Primenu odabrane operacije na unete operande na osnovu izabrane opcije izdvajamo u posebnu funkciju \texttt{calculate}. U njoj ćemo proveriti da li je opcija validna, primeniti operaciju i vratiti rezultat. Ukoliko se ne pošalje dobra vrednost argumenta \texttt{choice} funkcija bi, npr., mogla da baci izuzetak tipa \texttt{std::invalid\_argument} ili vrati vrednost \texttt{false}. Proveru možemo uraditi i sa \texttt{assert(choice\ \textgreater{}=\ 1\ \&\&\ choice\ \textless{}=\ 4);} iz \texttt{cassert} zaglavlja. Ukoliko se ne pošalje očekivana vrednost za argument \texttt{choice} program će biti prekinut. Ovo nam je način da nametnemo važenje preduslova za funkciju. Istu proveru bismo mogli da stavimo i u \texttt{printResults} funkciju. Dobili smo funkcije koje rade nezavisne poslove. Time su potencijalno upotrebljive i za neku dalju upotrebu ili pisanje kompleksnijeg kalkulatora. Sada i \texttt{main} funkcija može izgledati jednostavno kao niz poziva ovih funkcija.
